\documentclass[11pt,a4paper]{article}
% Prefer XeLaTeX or LuaLaTeX; falls back safely under pdfTeX.
\usepackage{iftex}
\ifXeTeX
  \usepackage{fontspec}
  \usepackage[spanish,shorthands=off]{babel}
  \setmainfont{Latin Modern Roman}
  \IfFontExistsTF{DejaVu Sans Mono}
    {\setmonofont{DejaVu Sans Mono}}
    {\IfFontExistsTF{Consolas}{\setmonofont{Consolas}}{\setmonofont{Latin Modern Mono}}}
\else\ifLuaTeX
  \usepackage{fontspec}
  \usepackage[spanish,shorthands=off]{babel}
  \setmainfont{Latin Modern Roman}
  \IfFontExistsTF{DejaVu Sans Mono}
    {\setmonofont{DejaVu Sans Mono}}
    {\IfFontExistsTF{Consolas}{\setmonofont{Consolas}}{\setmonofont{Latin Modern Mono}}}
\else
  \usepackage[T1]{fontenc}
  \usepackage[utf8]{inputenc}
  \DeclareUnicodeCharacter{2248}{\ensuremath{\approx}}
  \usepackage{lmodern}
  \usepackage[spanish,shorthands=off]{babel}
\fi\fi

\usepackage[margin=23mm]{geometry}
\usepackage{xcolor}
\usepackage{fvextra}
\usepackage[hidelinks]{hyperref}

\definecolor{codegray}{HTML}{526173}
\definecolor{codebg}{HTML}{EDF2F7}
\DefineVerbatimEnvironment{Codigo}{Verbatim}{
  fontsize=\small,breaklines,breakanywhere,numbers=left,
  numbersep=8pt,xleftmargin=1.7em,rulecolor=\color{codegray},
  frame=single,framerule=0.6pt,framesep=6pt,
  fillcolor=\color{codebg},tabsize=4
}

\setlength{\parindent}{0pt}
\setlength{\parskip}{5pt}

\title{Proyecto Final POO\\Rutas ciclistas sobre OpenStreetMap}
\author{Explicación técnica e historial de cambios}
\date{}

\begin{document}
\maketitle
\tableofcontents
\clearpage

\section*{Cómo leer este documento}
Los fragmentos son extractos literales del repositorio. La leyenda indica
el archivo, la versión y el rango de líneas original. Un extracto puede
empezar dentro de una función y conserva la sangría del archivo. El documento
incluye el código en su propio archivo \texttt{.tex}.

\section{Fase 2. Coordenadas separadas, área, zoom y progreso}

\textbf{Conceptos de esta fase.} Streamlit ejecuta de nuevo el código de la interfaz al usar sus controles. El estado de sesión guarda datos entre ejecuciones. CLI significa interfaz de línea de comandos.

La web usaba un campo común para el desplazamiento del mapa. A raíz del reporte, el formulario pasó a leer y guardar el valor del extremo elegido. Cada extremo tiene su propia clave.

\subsection{UI: origen y destino independientes}
Al fijar un punto se actualiza exclusivamente su clave, se invalida el resultado anterior y se centra el mapa.

\smallskip
{\footnotesize\texttt{\detokenize{src/app/ui/streamlit_app.py}} · versión actual, líneas 242--264.}\par
\begin{Codigo}[firstnumber=242]
    def _coordinate_form(self, mode: str, point_key: str) -> None:
        current: Coordinate | None = st.session_state[point_key]
        value = f"{current.lat},{current.lon}" if current is not None else ""
        name = mode.lower()
        with st.form(f"coordinates_{point_key}"):
            raw = st.text_input(
                f"Coordenadas de {name} (LAT,LON)",
                value=value,
                key=f"{point_key}_coordinates_{value}",
            )
            if st.form_submit_button(f"Fijar {name}"):
                try:
                    coordinate = Coordinate.from_string(raw)
                except ValueError as exc:
                    st.error(str(exc))
                else:
                    self._set_point(point_key, coordinate)
                    self._recenter(coordinate)

    @staticmethod
    def _set_point(point_key: str, coordinate: Coordinate) -> None:
        st.session_state[point_key] = coordinate
        st.session_state["route_result"] = None
\end{Codigo}

\textbf{Lectura del fragmento.} \texttt{self} es la interfaz; \texttt{mode} y \texttt{point\_key} son textos. \texttt{st.session\_state[clave]} guarda valores entre reejecuciones. \texttt{current} es una \texttt{Coordinate} o \texttt{None}. \texttt{value} es el texto inicial y \texttt{name} una copia de \texttt{mode} en minúsculas obtenida con \texttt{lower()}. Las f-strings interpolan valores y la expresión condicional produce una cadena vacía si el punto falta. \texttt{st.form(identificador)} crea el contexto del formulario. \texttt{st.text\_input(etiqueta, value=..., key=...)} devuelve \texttt{raw}, el texto introducido. \texttt{st.form\_submit\_button(etiqueta)} devuelve un booleano. \texttt{Coordinate.from\_string(texto)} devuelve \texttt{coordinate} o lanza \texttt{ValueError}. \texttt{st.error(str(exc))} muestra el mensaje y \texttt{str} convierte el error a texto. El \texttt{else} del \texttt{try} solo actúa si se pudo convertir. \texttt{\_set\_point(clave, coordenada)} cambia ese extremo e invalida \texttt{route\_result}. \texttt{@staticmethod} indica que la función no recibe \texttt{self} automáticamente. \texttt{\_recenter(coordenada)} actualiza la posición del mapa. Ambos métodos devuelven \texttt{None}.


\subsection{Configuración y mapa: área automática y zoom estable}
La web y la terminal amplían el área al tamaño requerido por los puntos y el margen. La CLI admite un límite manual explícito. Para evitar que el mapa se pusiera blanco y regresara a su posición inicial, el centro y el zoom iniciales del componente se mantienen estables mientras se guardan los cambios informados por Leaflet.

\smallskip
{\footnotesize\texttt{\detokenize{src/app/config.py}} · versión actual, líneas 92--97.}\par
\begin{Codigo}[firstnumber=92]
    def with_area_for(self, origin: Coordinate, dest: Coordinate) -> Settings:
        """Amplía el límite base hasta cubrir los puntos y el margen."""
        area = BoundingBox.from_points(origin, dest, self.margin_m).area_km2
        if area <= self.max_area_km2:
            return self
        return replace(self, max_area_km2=area)
\end{Codigo}

\textbf{Lectura del fragmento.} \texttt{origin} y \texttt{dest} son coordenadas y \texttt{self} representa la configuración \texttt{Settings}. \texttt{BoundingBox.from\_points(origen, destino, margen\_en\_metros)} devuelve el rectángulo. Su atributo \texttt{area\_km2} da \texttt{area}, un número decimal. \texttt{\textless{}=} comprueba si cabe en el límite. \texttt{replace} pertenece a \texttt{dataclasses}, biblioteca estándar de Python: \texttt{replace(objeto, campo=valor)} devuelve una copia con el campo modificado, apropiada cuando las opciones se tratan como valores inmutables.


\smallskip
{\footnotesize\texttt{\detokenize{src/app/ui/streamlit_app.py}} · versión actual, líneas 171--194.}\par
\begin{Codigo}[firstnumber=171]
        if st.session_state.get("map_points") != (origin, dest):
            st.session_state["map_points"] = (origin, dest)
            st.session_state["map_initial_view"] = view
        initial_view = st.session_state["map_initial_view"]

        map_obj = self._renderer.render(
            origin=origin,
            dest=dest,
            center=initial_view.center,
            zoom=initial_view.zoom,
        )
        map_output = st_folium(
            map_obj,
            height=480,
            key=f"click_map_{st.session_state.get('map_revision', 0)}",
            returned_objects=["last_clicked", "center", "zoom"],
        )
        if isinstance(map_output, dict) and map_output.get("center"):
            center = map_output["center"]
            with suppress(ValueError, TypeError, KeyError):
                self._save_view(
                    MapView(Coordinate(center["lat"], center["lng"]), map_output["zoom"])
                )
        clicked = map_output.get("last_clicked") if isinstance(map_output, dict) else None
\end{Codigo}

\textbf{Lectura del fragmento.} \texttt{origin} y \texttt{dest} son coordenadas opcionales. \texttt{view} e \texttt{initial\_view} son objetos \texttt{MapView} con centro y zoom. \texttt{st.session\_state.get(clave, predeterminado)} lee estado persistente durante la sesión. \texttt{!=} compara la tupla de puntos con la anterior. \texttt{\_renderer.render(origin=..., dest=..., center=..., zoom=...)} devuelve \texttt{map\_obj}, un mapa Folium. \texttt{st\_folium(mapa, height=..., key=..., returned\_objects=...)} integra el mapa y devuelve \texttt{map\_output}, un diccionario de eventos. \texttt{key} identifica el componente; se construye con una revisión entera. \texttt{isinstance(valor, dict)} valida el tipo y \texttt{get} permite leer campos opcionales. \texttt{center} es el diccionario con latitud y longitud; \texttt{clicked} es otro diccionario o \texttt{None}. \texttt{suppress(tipos\_de\_error)} viene de \texttt{contextlib} y omite los errores enumerados dentro del \texttt{with}. \texttt{Coordinate(latitud,longitud)} y \texttt{MapView(coordenada,zoom)} construyen los valores; \texttt{\_save\_view(vista)} los persiste si cambian.


\subsection{Terminal y web: progreso, PNG y resumen}
La barra indica etapas de preparación, cálculo y generación de imagen. El proveedor carece de un porcentaje de descarga que pueda mostrarse. \texttt{--no-summary} oculta métricas e instrucciones por calle y mantiene progreso y avisos. \texttt{--no-show} guarda el PNG sin ventana. \texttt{--no-plot} omite el PNG. Para pares negativos se usa, por ejemplo, \texttt{--origin=-12.01972,-77.04861}.

\smallskip
{\footnotesize\texttt{\detokenize{src/app/ui/cli.py}} · versión actual, líneas 78--85.}\par
\begin{Codigo}[firstnumber=78]
            self._progress(0, "Preparando red ciclista (caché o descarga)…")
            planner.prefetch(origin, dest)
            self._progress(1, "Red lista. Calculando ruta…")
            result = planner.plan_with_alternatives(origin, dest)
        except (RoutingError, ValueError) as exc:
            print(f"Error: {exc}", file=sys.stderr)
            if isinstance(exc, RoutingError) and exc.diagnostic:
                print(f"Diagnóstico\n{exc.diagnostic}", file=sys.stderr)
\end{Codigo}

\textbf{Lectura del fragmento.} \texttt{planner} es un planificador y \texttt{origin}, \texttt{dest} son coordenadas. \texttt{\_progress(etapa\_entera, mensaje)} imprime el estado; \texttt{planner.prefetch(origen,destino)} prepara la red y devuelve \texttt{None}; \texttt{plan\_with\_alternatives(origen,destino)} devuelve \texttt{result}, un \texttt{RouteResult}. \texttt{except (RoutingError, ValueError) as exc} captura cualquiera de esos errores. \texttt{sys.stderr} es el canal de errores de la terminal. \texttt{print(texto, file=canal)} escribe allí; \texttt{isinstance} comprueba si el error admite un \texttt{diagnostic}, un mensaje opcional. La secuencia \texttt{\textbackslash{}n} representa un salto de línea.


\smallskip
{\footnotesize\texttt{\detokenize{src/app/ui/cli.py}} · versión actual, líneas 95--121.}\par
\begin{Codigo}[firstnumber=95]
        if not args.no_summary:
            self._print_result(result)
            self._print_instructions(result)
        elif result.alternative_origin is not None or result.alternative_dest is not None:
            for warning in result.warnings:
                print(f"Aviso: {warning}", file=sys.stderr)
        # Generar visualización si está activada
        if getattr(args, "plot", True):
            self._progress(2, "Ruta calculada. Generando el mapa…")
            try:
                plotter = RoutePlotter(
                    margin_m=args.margin_m,
                    dpi=args.dpi,
                    labels=not args.no_labels,
                    basemap=None if args.basemap == "none" else args.basemap,
                )
                out_path = getattr(args, "plot_out", "route.png") or "route.png"
                no_show = getattr(args, "no_show", False)
                plotter.render(result, out_path=out_path, show=not no_show)
                for note in plotter.warnings:
                    print(f"Aviso: {note}", file=sys.stderr)
            except Exception as exc:
                print(f"Aviso: no se pudo generar la imagen: {exc}", file=sys.stderr)
                return 0
            self._progress(3, "Ruta y mapa listos.")
        else:
            self._progress(3, "Ruta calculada.")
\end{Codigo}

\textbf{Lectura del fragmento.} \texttt{args} guarda opciones de la CLI y \texttt{result} es el resultado de ruta. \texttt{\_print\_result(resultado)} y \texttt{\_print\_instructions(resultado)} escriben resumen e indicaciones. \texttt{elif} evalúa otra condición si la primera falló. \texttt{warning} y \texttt{note} son textos de aviso. \texttt{getattr(objeto, nombre, predeterminado)} lee un atributo o usa el valor de respaldo. \texttt{RoutePlotter(margin\_m=..., dpi=..., labels=..., basemap=...)} crea \texttt{plotter}, el renderizador; margen y resolución son números, etiquetas es booleano y fondo es texto o \texttt{None}. \texttt{out\_path} es la ruta de salida y \texttt{no\_show} un booleano; \texttt{not} lo invierte para activar \texttt{show}. \texttt{plotter.render(resultado, out\_path=..., show=...)} escribe el PNG y devuelve su ruta. \texttt{plotter.warnings} es la colección de avisos. La excepción \texttt{exc} preserva el fallo del PNG y \texttt{return 0} indica que el cálculo de la ruta ya había tenido éxito.


\smallskip
{\footnotesize\texttt{\detokenize{src/app/ui/streamlit_app.py}} · versión actual, líneas 52--66.}\par
\begin{Codigo}[firstnumber=52]
def _prefetch_progress(future: Future[None], calculate_when_ready: bool = False) -> None:
    if not future.done():
        message = "Preparando red ciclista (caché o descarga)…"
        if calculate_when_ready:
            message += " La ruta se calculará al terminar."
        st.progress(0, text=message)
    elif calculate_when_ready:
        st.rerun()
    elif not future.cancelled() and future.exception() is None:
        st.progress(50, text="Red lista. Pulsa Calcular ruta.")
    else:
        st.warning("No se pudo preparar la red. Pulsa Calcular ruta para ver el detalle.")


class StreamlitApp:
\end{Codigo}

\textbf{Lectura del fragmento.} \texttt{Future[None]} representa una tarea que podrá terminar más tarde y cuyo trabajo devuelve \texttt{None}. \texttt{future} representa la preparación del grafo. \texttt{calculate\_when\_ready} es booleano con valor inicial \texttt{False}. \texttt{future.done()} y \texttt{future.cancelled()} devuelven booleanos. \texttt{future.exception()} devuelve el error o \texttt{None} cuando la tarea terminó. \texttt{message} es texto y \texttt{+=} concatena otra frase. \texttt{st.progress(entero, text=...)} muestra una etapa y devuelve el componente que se puede actualizar. \texttt{st.rerun()} solicita una nueva ejecución de la interfaz. \texttt{st.warning(texto)} muestra un aviso. La rama de progreso distingue una preparación pendiente de una ya disponible.


\end{document}
